\documentclass[11pt]{article}

\usepackage[margin=1in]{geometry}
\usepackage{amsmath,amssymb}

\title{A Comparison of Two Ratios}
\author{}
\date{}

\begin{document}

\maketitle

\section*{Statement}

Let \(A>B>C\), let \(b\in(0,1)\) be fixed, and let \(d,e>0\). Suppose both denominators are positive:
\[
B-bC>0
\qquad\text{and}\qquad
B-e-b(C-d)>0.
\]
Then
\[
\frac{A-bC}{B-bC}
<
\frac{A-e-b(C-d)}{B-e-b(C-d)}
\]
if and only if
\[
e>bd.
\]

\section*{Proof}

Define
\[
X=A-bC,\qquad Y=B-bC,\qquad t=bd-e.
\]
Then
\[
A-e-b(C-d)
=A-e-bC+bd
=X+t,
\]
and similarly
\[
B-e-b(C-d)
=B-e-bC+bd
=Y+t.
\]
Thus the desired inequality is
\[
\frac{X}{Y}<\frac{X+t}{Y+t}.
\]
Since the two denominators are positive by assumption, we may subtract the left-hand side from the right-hand side without changing signs in the denominator:
\[
\frac{X+t}{Y+t}-\frac{X}{Y}
=\frac{Y(X+t)-X(Y+t)}{Y(Y+t)}
=\frac{t(Y-X)}{Y(Y+t)}.
\]
Now
\[
Y-X=(B-bC)-(A-bC)=B-A<0,
\]
because \(A>B\). Also \(Y>0\) and \(Y+t>0\) by assumption. Therefore
\[
\frac{X+t}{Y+t}-\frac{X}{Y}>0
\]
if and only if
\[
t<0.
\]
Since \(t=bd-e\), this is equivalent to
\[
bd-e<0,
\]
or
\[
e>bd.
\]
Hence
\[
\frac{A-bC}{B-bC}
<
\frac{A-e-b(C-d)}{B-e-b(C-d)}
\]
if and only if \(e>bd\), as claimed.

\end{document}
